\documentclass[10pt]{article}

\usepackage{fancyhdr}
\pagestyle{fancy}
\fancyhf{}
\fancyhead[C]{Working Note: Tightness and Evaluation of the Learner's Requisite Variety}
\fancyfoot[C]{\thepage}
\renewcommand{\headrulewidth}{0.4pt}

\usepackage[margin=2cm]{geometry}
\usepackage[T1]{fontenc}
\usepackage[utf8]{inputenc}
\usepackage{amsmath,amssymb,mathtools}
\usepackage{booktabs}
\usepackage{enumitem}

\setlist{nosep,leftmargin=1.3em}
\newcommand{\E}{\mathbb E}

\begin{document}

\noindent\textbf{Parked note, 2026-11-26.}
Written to revisit much later.
It records what is known about how tight Corollary 1 of the RLC paper is (\texttt{notes\_joery/rlc/sections/split.tex}, Eq.~6 in the build of this date) and how $R(\delta)$ could be evaluated.
Nothing here is in the paper.

\section{The statement}
With $W_t = f(\hat\theta_t, \hat S_t)$ a hyperstate statistic that determines the within-tolerance regulators $\Pi^\circ$, a fixed point $\pi^\star$ with $\Pr(\pi^\star \notin \Pi^\circ(W_\tau)) \le \delta$ requires
\begin{align*}
  \E[\tau]\, \big(I_{\hat\theta} + I_{\hat S}\big) \;&\ge\; R(\delta) + \E[\tau]\, \bar\Delta - \big(H(W_\tau) - H(W_0)\big), \\
  R(\delta) \;&=\; \inf\nolimits_{p(\pi \mid W_\tau)} \big\{ I(W_\tau; \pi) \,:\, \Pr\big(\pi \notin \Pi^\circ(W_\tau)\big) \le \delta \big\}.
\end{align*}
The bound is necessary and not sufficient.

\section{Is it tight?}
It is sharp, since equality is attained in some models, and it is not tight for a typical learner.
The proof chains three inequalities, each with its own slack.
\begin{enumerate}
  \item \textbf{Proposition 4.}
  $\E[\tau](I_{\hat\theta} + I_{\hat S}) \ge H(W_0) - H(W_\tau \mid X_\tau) + \E[\tau]\bar\Delta$, with equality when $f$ is the identity.
  For a smaller $W_t$ the left side also counts information about the rest of the hyperstate.
  \item \textbf{History to policy.}
  $I(W_\tau; X_\tau) \ge I(W_\tau; \pi^\star)$ by data processing.
  Equality needs the history to say nothing about $W_\tau$ beyond what the choice of $\pi^\star$ reveals.
  \item \textbf{Policy to optimum.}
  $I(W_\tau; \pi^\star) \ge R(\delta)$, since $\pi^\star$ is one feasible channel.
  Equality needs the learner's choice to be the least informative one that reaches error $\delta$.
\end{enumerate}
\textbf{Equality case.}
Two equally likely stationary hypotheses, each with its own single regulator, $\delta = 0$, $\tau = 1$, and one outcome that reveals the hypothesis.
Both sides equal $\ln 2$.

\noindent\textbf{Where the slack comes from.}
Mostly steps 1 and 2, a learner learns things that the decision does not need.
This is the same gap as on the sufficient side, where a floor on $I_{\hat\theta} + I_{\hat S}$ cannot force its information to be about $W_t$.

\section{Evaluating $R(\delta)$}
$R(\delta)$ is a rate--distortion function with the 0/1 distortion $\mathbf 1\{\pi \notin \Pi^\circ(w)\}$.
It depends only on the learner's model, the law of $W_\tau$ and the tolerance $\varepsilon_\pi$.
\begin{itemize}
  \item \textbf{Properties.}
  $0 \le R(\delta) \le R(0) \le H(W_\tau)$ for discrete $W_\tau$, non-increasing and convex in $\delta$, zero from $\delta = 1 - p_{\max}$ on, with $p_{\max} = \sup_\pi \Pr(\pi \in \Pi^\circ(W_\tau))$ the largest belief mass that one policy regulates.
  It is an infimum in general and a minimum for finite $W_\tau$ and finite policy sets.
  \item \textbf{Exact.}
  Blahut--Arimoto on the membership table of policies against hyperstates.
  \item \textbf{Sampled.}
  Draw $w_1, \dots, w_n$ from $\hat b_0$, solve each for a regulator, evaluate every regulator on every sample, and run Blahut--Arimoto on the $n \times n$ table.
  \item \textbf{Upper bound.}
  $R(\delta) \le H(\pi(W_\tau))$ for any assignment $w \mapsto \pi(w) \in \Pi^\circ(w)$.
  \item \textbf{Lower bound.}
  $R(\delta) \ge d(1 - \delta \,\|\, p_{\max})$ for $1 - \delta \ge p_{\max}$, with $d$ the binary KL divergence, by data processing of the KL between the joint law of $(W_\tau, \pi)$ and the product of its marginals on the event $\pi \in \Pi^\circ(W_\tau)$.
  \item \textbf{Zero error.}
  $R(0) = \inf_m \E\big[-\ln m(\Pi^\circ(W_\tau))\big]$ over priors $m$ on policies, with the optimal channel $p(\pi \mid w) \propto m(\pi)\, \mathbf 1\{\pi \in \Pi^\circ(w)\}$.
  For $\delta > 0$ the indicator softens to $\exp(-\beta\, \mathbf 1\{\pi \notin \Pi^\circ(w)\})$.
\end{itemize}

\section{Verification status}
\begin{center}
\begin{tabular}{p{0.52\linewidth}p{0.4\linewidth}}
\toprule
Claim & Status \\
\midrule
The inequality of Section 1, stationary and drifting $W_t$ & Checked on 500 random finite models each, no violations. \\
Continuous hypothesis class & Checked by quadrature on 12 models. \\
Equality case & By hand. \\
Typical size of the gap & Not measured. \\
Lower bound through $p_{\max}$ & Standard argument, checked only on the three-hypothesis tie example ($\log_2 \tfrac32 = 0.585$ bits, exact). \\
Formula for $R(0)$ & Standard, checked only on the same example. \\
Zero difference $H(W_\tau) - H(W_0)$ for a stationary latent state with a history-dependent stopping time & Believed false in general, not run. \\
\bottomrule
\end{tabular}
\end{center}

\section{To revisit}
\begin{enumerate}
  \item Measure the gap of the three steps separately on random models, to see which dominates.
  \item A floor on $I_W$ in place of $I_{\hat\theta} + I_{\hat S}$ removes step 1 and is what the sufficient side needs, estimated by posterior sampling.
  \item Verify the $p_{\max}$ bound and the $R(0)$ formula numerically, then decide whether $p_{\max}$ belongs in the text as the practical scale for $c_{\hat\theta}$.
  \item Whether a learner can be built that attains step 3, a policy choice that is the rate--distortion optimum, and whether that relates to posterior sampling over regulators.
\end{enumerate}

\end{document}
